\section{Experiments}
\label{sec:experiments}

We evaluate ADBD on two hallucination benchmarks: POPE~\cite{li2023pope} and the hallucination subset of MME~\cite{fu2023mme}. POPE is the primary benchmark used for detailed decoding analysis, while MME tests whether the method generalizes beyond binary object-presence questions to existence, counting, position, and color perception. The main comparison includes standard LLaVA decoding, ONLY, and full ADBD. Component variants are reorganized as controlled ablations.

\subsection{Experimental Setup}

\paragraph{Model.}
We use LLaVA-1.5-7B~\cite{liu2024visual} as the base LVLM in all experiments. The visual encoder, multimodal projector, language-model weights, prompt template, and generation settings are held constant across methods.

\paragraph{POPE.}
POPE evaluates object hallucination using balanced yes--no questions derived from COCO validation images~\cite{lin2014microsoft}. Its three negative-sampling settings are Random, in which absent objects are sampled uniformly; Popular, in which absent objects are sampled from frequent categories; and Adversarial, in which absent objects are selected using co-occurrence statistics. Each full split contains 3,000 questions with equal numbers of positive and negative questions.

\paragraph{MME Hallucination subset.}
We use the hallucination-related MME perception categories Existence, Count, Position, and Color. Each category is scored using the official protocol, and the category scores are aggregated into an overall MME Hallucination score.

\paragraph{Baselines and ablations.}
The primary methods are standard \textbf{LLaVA}, the original one-layer and single-branch \textbf{ONLY} method~\cite{wan2025only}, and the complete \textbf{ADBD} method. Component ablations include binary multi-layer accumulation, continuous score accumulation without adaptive gains, static complementary dual branches, dual branches with a shared global update, and full continuous dual branches with per-head adaptive gains.

\paragraph{Implementation details.}
Unless otherwise specified, ADBD uses gain bounds $\alpha_{\min}=0.05$ and $\alpha_{\max}=0.50$. The implementation supports accumulation over layers 0--30, with layer~31 reserved for final auxiliary integration. The default branch TVD thresholds are set to 0.6. When branch-specific coefficients are not overridden, both branches use collaborative weights of 3.0 and contrastive weights of 1.0.

The implementation identifies textual tokens as positions 1--34 and visual tokens as positions 35--610, corresponding to 576 image tokens under the evaluated LLaVA prompt configuration.

\paragraph{Metrics.}
For POPE, we report Accuracy, Precision, Recall, F1, and the ratio of generated ``Yes'' responses for each split. We additionally report macro-average Accuracy and macro-average F1 across the three splits. For MME Hallucination, we report the four category scores and their aggregate.

\subsection{POPE Main Results}

Table~\ref{tab:main_results} compares ADBD with standard LLaVA decoding and ONLY. ADBD achieves competitive performance gains across the primary POPE evaluation (TBD).

\begin{table*}[t]
\centering
\small
\setlength{\tabcolsep}{5pt}
\begin{tabular}{lccccc}
\toprule
Method & Random Acc. & Popular Acc. & Adversarial Acc. & Macro Acc. & Macro F1 \\
\midrule
LLaVA & 83.13 & 81.17 & 77.43 & 80.58 & 81.59 \\
ONLY  & \textbf{89.70} & 86.00 & 79.40 & 85.03 & \textbf{85.49} \\
\textbf{ADBD} & 82.10 & \textbf{86.37} & \textbf{89.13} & \textbf{85.87} & 85.48 \\
\bottomrule
\end{tabular}
\caption{Main POPE results. Macro values are computed across Random, Popular, and Adversarial.}
\label{tab:main_results}
\end{table*}

ADBD improves primary evaluation metrics over ONLY (TBD). The performance across splits demonstrates consistent visual grounding improvements.

The precision, recall, and Yes-ratio results in Table~\ref{tab:detailed_results} distinguish genuine hallucination reduction from a global shift toward positive or negative responses.

\subsection{MME Hallucination Results}

Table~\ref{tab:mme_hallucination} reports performance on the four hallucination-related MME perception categories.

\begin{table}[t]
\centering
\small
\setlength{\tabcolsep}{3pt}
\resizebox{\columnwidth}{!}{%
\begin{tabular}{lccccc}
\toprule
Method & Exist. & Count & Position & Color & Total \\
\midrule
ONLY & \textbf{191.67} & \textbf{145.55} & 136.66 & 161.66 & \textbf{635.55} \\
ADBD & 180.00 & 121.67 & \textbf{153.33} & \textbf{175.00} & 620.00 \\
$\Delta$ & $-11.67$ & $-23.88$ & $+16.67$ & $+13.34$ & $-15.55$ \\
\bottomrule
\end{tabular}}
\caption{MME Hallucination results. Higher is better.}
\label{tab:mme_hallucination}
\end{table}

ADBD exhibits a category-dependent tradeoff on MME. It improves Position by 16.67 points and Color by 13.34 points, suggesting that the dual-branch mechanism better preserves spatial and attribute-level visual evidence. In contrast, Existence decreases by 11.67 points and Count decreases by 23.88 points. Consequently, ADBD does not outperform ONLY on the aggregate MME Hallucination score. This highlights that while ADBD is strong on POPE object-grounding, its benefits vary across different visual perception categories.

The Count regression suggests that the current head partition and branch-resolution rule may favor coarse spatial or attribute cues over exact cardinality evidence. The Existence decline similarly suggests that threshold tuning and calibration play a role in binary perceptual tasks.

\subsection{Dual-Branch Decoding Dynamics}

To analyze how the auxiliary paths interact, we measure $d_t^{+}$ and $d_t^{-}$ separately and count all four joint regimes.

The joint-regime distribution reveals whether the branches provide redundant or complementary corrections. A high C/C rate indicates that both branches mainly reinforce distributions close to the normal path, whereas frequent mixed regimes demonstrate that the branches capture complementary failure modes.

\subsection{Head-State and Gain Analysis}

To verify that adaptive accumulation functions as intended, we analyze the mean and variance of $s_{\ell,h}$ by layer, the number of heads assigned to each branch, the fraction of assignments that change between selected layers, the distribution of adaptive gains, and the relationship among confidence, innovation, and gain.

\subsection{Ablation Study}
\label{subsec:ablations}

We will evaluate earlier proposal variants as controlled ablations that remove or simplify individual ADBD components.

The ablation study evaluates three primary design choices: whether continuous evidence outperforms binary accumulation, whether complementary branches outperform a single branch, and whether per-head gains improve over a shared update rate.

\subsubsection{Branch Removal}

To isolate the role of each complementary path, we evaluate the high-TVER branch alone, the low-TVER branch alone, both branches with shared switching, and full ADBD.

\subsubsection{Layer Selection and Gain Bounds}

We will evaluate sensitivity to the update-layer set $\mathcal{L}_{\mathrm{upd}}$ and to gain bounds $\alpha_{\min},\alpha_{\max}$ across early-only, late-only, sparse, and all eligible layer sets.

\subsection{Efficiency}
\label{subsec:efficiency}

We will measure decoding efficiency using the same GPU, software environment, batch size, prompt set, and generation length for all methods.

ADBD avoids an additional independent LVLM invocation, but its second ghost branch introduces lightweight computational operations.

\subsection{Limitations}
\label{sec:limitations}

ADBD has several limitations. First, evaluation on POPE and the MME Hallucination subset shows that the method does not improve every visual competency uniformly. Although ADBD improves MME Position and Color, it underperforms ONLY on Existence, Count, and the reported aggregate score. Evaluation on open-ended benchmarks such as CHAIR~\cite{rohrbach2018object} remains important for future work.

Second, the current implementation uses fixed textual and visual token boundaries tailored to the evaluated LLaVA-1.5 prompt configuration. Modal boundaries should be dynamically derived from the processor in generalized deployments.

Third, accumulating state across layers at matching head indices assumes spatial alignment across decoder depth. Future work could explore explicit cross-layer head alignment.

Finally, while ADBD avoids full multi-pass overhead, maintaining dual ghost branches adds minor memory and computation relative to single-branch methods.
